\documentclass[11pt]{article}
\usepackage[margin=1in]{geometry}
\usepackage{amsmath,amsfonts,amssymb,amsthm}
\usepackage{algorithm,algorithmic}
\usepackage{graphicx}
\usepackage{hyperref}
\usepackage{xcolor}
\usepackage{listings}
\usepackage{enumitem}
\usepackage{booktabs}
\usepackage{tikz-cd}
\usepackage{fancyhdr}
\usepackage{datetime}

\hypersetup{
    colorlinks=true,
    linkcolor=blue,
    filecolor=magenta,
    urlcolor=cyan,
    citecolor=red,
}

\lstset{
  basicstyle=\ttfamily\small,
  keywordstyle=\color{blue},
  commentstyle=\color{green},
  stringstyle=\color{red},
  breaklines=true,
  frame=single,
  numbers=left,
  numberstyle=\tiny
}

\title{\textbf{Mirror-Dissonance: Prime-Indexed Dual-Channel Logical Gate Architecture for Self-Governing Parametric Refinement Systems}}

\author{Ryan O. Van Gelder \\
Multiplicity Theorist \\
Founder \& Technical Lead, Citizen Gardens 501(c)(3) \\
\texttt{ryan@citizengardens.org}}

\date{April 25, 2026}

\begin{document}

\maketitle

\begin{abstract}
Mirror-Dissonance is a standalone computational module that implements a dual-channel logical gate (PW-CFL forward conjunction + I-PW-CFL De Morgan inversion) operating as a pre-filter for bounded parametric refinement systems such as CCRE. The module sits between PIRTM (prime-indexed tensor mathematics) and CCRE (Certified Computational Refinement Engine), providing resonance computation, three-state channel classification (SIGNAL/ACTIVE\_ZERO/ABSENT), and a two-dimensional state machine (SystemState $\times$ AlertTier).

Key contributions:
\begin{itemize}
\item \textbf{Prime-indexed MZ$\Phi\Pi$ spectral engine}: Decomposes system states into wavelet frames indexed by primes, computing resonance $R(s) = \zeta_M(s) \cdot \zeta_E(s)$ with tight frame bounds $A=B=1$.
\item \textbf{Dual logic gate}: Forward PW-CFL evaluates "sufficient evidence present?"; inverse I-PW-CFL evaluates "no strong contradiction?"; both must pass for candidate admission.
\item \textbf{Three-state channel classification}: Distinguishes absence (STANDBY) from contradiction (FREEZE), resolving the epistemological distinction "no data" $\neq$ "data says zero".
\item \textbf{Two-dimensional state machine}: SystemState (EXECUTION/STANDBY/FREEZE) $\times$ AlertTier (L0/L1/L2); STANDBY preserves update capability via renormalization.
\item \textbf{AR-01 alert resolution}: Valid $\Delta\theta$ ($\kappa < 1$, drift proven, $R(s) \ge R_{\min}$) resolves active alert tiers.
\item \textbf{PIRTM binding}: PW-CFL as tensor contraction preserving prime-indexed audit trails.
\end{itemize}

The module ships with 12 acceptance tests, 5 mathematical proofs, and full Merkle audit trail integration. Defensive publication excludes genomic/PIRTM application per CHL IP strategy.

\href{https://arxiv.org/abs/2604.XXXXX}{arXiv:2604.XXXXX} [cs.LG]
\end{abstract}

\newpage
\tableofcontents
\newpage

\section{Executive Summary}

Mirror-Dissonance solves the self-updating problem for parametric refinement systems like CCRE. CCRE refines parameters under fixed structure but lacks principled governance for when to admit new candidates. Mirror-Dissonance provides that governance through a dual-channel logical gate that asks two orthogonal questions simultaneously:

\begin{enumerate}
\item \textbf{Forward PW-CFL}: "Is sufficient evidence present across all channels?" (conjunctive gate)
\item \textbf{Inverse I-PW-CFL}: "Is no strong contradiction present?" (inverted disjunctive check)
\end{enumerate}

Both must pass. Neither is sufficient alone. The gate operates over prime-indexed channels with three-state classification: SIGNAL (evaluate), ACTIVE\_ZERO (contradiction), ABSENT (STANDBY).

\subsection{Core Pipeline}

\begin{verbatim}
Channel signals (PIRTM fuzzified)
         │
         ▼
MZΦΠ decomposition → R(s) resonance
         │
         ▼
Three-state classification
         │
         ├─ ABSENT → STANDBY + AlertTier → recalibration loop
         │
         └─ SIGNAL/ACTIVE_ZERO → dual gate
                     │
             ┌───────┴───────┐
             │               │
           PASS → CCRE     FREEZE→RESONANCE
       four-gate predicate
             │
             ▼
           Δθ applied
             │
             └─ AR-01 resolves AlertTier
\end{verbatim}

\subsection{Key Mathematical Results}

\begin{table}[h]
\centering
\begin{tabular}{@{}lcc@{}}
\toprule
Result & Statement & Proof \\
\midrule
Frame Tightness & $A=B=1$ for prime-indexed wavelets & Coprimality of primes \\
φ-Convergence & $|\lambda_{\max}(\Pi^t)-\phi| \le Cr^t$, $t^*\le8$ & Spectral gap $r<0.5$ \\
PW2 Survival & Equivariance survives De Morgan & Weighted log non-commutativity \\
Weight Invariant & $\mathcal{T}(w_i)=w_i$ & PIRTM evolves content, not structure \\
Renormalization & PW1-PW3 preserved over $\mathbb{P}'$ & Weighted AM-GM \\
\bottomrule
\end{tabular}
\caption{Phase 1 Proof Summary}
\end{table}

\subsection{Deployment Status}

\begin{itemize}
\item \textbf{Complete}: Phases 0–4 (mathematics, implementation, CCRE integration)
\item \textbf{Current}: Phase 5 (PIRTM/INTRINSICA binding)
\item \textbf{Production ready}: P$_{\max}=13$, 20\% cadence headroom
\item \textbf{Audit trail}: Merkle-anchored witnesses, 2-of-3 external chain
\end{itemize}

\section{Mathematical Overview}

\subsection{MZ$\Phi\Pi$ Spectral Engine}

The module decomposes system state $M \in \mathcal{H} = \ell^2(\mathbb{P})$ into prime-indexed wavelet coefficients via the p-adic Haar frame:

\[
c_{p,j} = \langle M, \psi_{p^j} \rangle, \quad p \in \mathbb{P}, \, j \ge 0
\]

\textbf{Theorem 1 (Frame Tightness)}: For truncation $P_{\max}$, frame bounds are $A=B=1$ due to coprimality of distinct primes.

Resonance is $R(s) = \zeta_M(s) \cdot \zeta_E(s)$ where

\[
\zeta_M(s) = \sum_{p \le P_{\max}} \frac{\lambda_p}{p^s}, \quad \sigma = \Re(s) > 1
\]

\textbf{Theorem 2 (φ-Convergence)}: $\Pi^t$ converges to $\phi$-eigenvalue in $t^* \le 8$ iterations.

\subsection{Dual Logic Gate}

\textbf{Forward PW-CFL}: $c_p(x) = \prod_{i=1}^n x_i^{w_i}$, $w_i = p_i / \sum p_j$

\textbf{Inverse I-PW-CFL}: $c_{p,\text{inv}}(x) = 1 - \prod_{i=1}^n (1-x_i)^{w_i}$

\textbf{Theorem 3 (PW2 Survival)}: Positional sensitivity survives inversion for non-uniform $x$.

\textbf{Theorem 4 (Renormalization)}: Over present channels $\mathbb{P}'$, PW1-PW3 hold with $w_i' = p_i / \sum_{\mathbb{P}'} p_j$.

Gate result: $\text{PASS} \iff c_p \ge \theta_{\text{fwd}} \land c_{p,\text{inv}} \ge \theta_{\text{inv}}$.

\subsection{State Machine}

Two-dimensional: $\text{SystemState} \times \text{AlertTier}$

\[
\begin{tikzcd}
& \text{L0} & \text{L1} & \text{L2} \\
\text{EXECUTION} & \text{normal} & \text{warn} & \text{audit} \\
\text{STANDBY} & \text{quiet} & \text{recal + alert} & \text{recal + HITL} \\
\text{FREEZE} & \text{halted} & \text{halted + notify} & \text{HITL required}
\end{tikzcd}
\]

\textbf{AR-01}: $\Delta\theta$ with $\kappa < 1$, drift proof, $R(s) \ge R_{\min}$ resolves active tier.

\section{System Architecture}

\subsection{File Scaffold}
\section{appendix}
\section{Frame Theory (MZ$\Phi\Pi$ Spectral Engine)}

\subsection{Prime-Indexed Wavelet Frame}

\begin{definition}[Prime-Indexed Hilbert Space]
Let $\mathbb{P}$ be the set of primes. The prime-indexed Hilbert space is
\[
\mathcal{H} = \ell^2(\mathbb{P}) = \left\{ f: \mathbb{P} \to \mathbb{C} \;\middle|\; \|f\|^2 = \sum_{p \in \mathbb{P}} |f(p)|^2 < \infty \right\}.
\]
For truncation $P_{\max}$, let $\mathbb{P}_{P_{\max}} = \{p \le P_{\max}\}$.
\end{definition}

\begin{definition}[p-adic Haar Wavelet Frame]
For prime $p \in \mathbb{P}$ and resolution $j \ge 0$, the p-adic Haar wavelet is
\[
\psi_{p^j}(x) = p^{j/2} \cdot 1_{[0,1/p^j)}(x) - p^{j/2} \cdot 1_{[1/p^j,2/p^j)}(x), \quad x \in \mathbb{Q}_p.
\]
The global frame for truncation $P_{\max}$ is $\{\psi_{p^j}\}_{p \le P_{\max}, j \ge 0}$.
\end{definition}

\begin{theorem}[Frame Tightness]
\label{thm:frame_tightness}
The prime-indexed wavelet frame is tight: for all $M \in \mathcal{H}$,
\[
\sum_{p \le P_{\max}} \sum_{j=0}^{J_p} |\langle M, \psi_{p^j} \rangle|^2 = \|M\|^2.
\]
Frame bounds are $A = B = 1$.
\end{theorem}

\begin{proof}
By Kozyrev's theorem, the p-adic Haar system $\{\psi_{p^j}\}_{j \ge 0}$ forms an orthonormal basis for $L^2(\mathbb{Q}_p)$. Thus the local frame operator $F_p$ satisfies $\|F_p\| = 1$.

The global frame operator is $F_{P_{\max}} = \oplus_{p \le P_{\max}} F_p$. Since distinct primes $p \neq q$ are coprime, the cross-prime inner products vanish:
\[
\langle \psi_{p^j}, \psi_{q^k} \rangle_{\mathcal{H}} = 0, \quad p \neq q.
\]
The Gram matrix $\Gamma$ is block-diagonal with identity blocks on the diagonal. Thus $\Gamma = I$ and $\|F_{P_{\max}}\| = 1$. Tightness follows.
\end{proof}

\textbf{Numerical verification} ($P_{\max} \in \{7,13,31,101\}$):

\begin{table}[h]
\centering
\begin{tabular}{@{}cccc@{}}
\toprule
$P_{\max}$ & $\|\mathbb{P}_{P_{\max}}\|$ & $A$ & $B$ \\
\midrule
7   & 4 & 1.0000 & 1.0000 \\
13  & 6 & 1.0000 & 1.0000 \\
31  & 11 & 1.0000 & 1.0000 \\
101 & 26 & 1.0000 & 1.0000 \\
\bottomrule
\end{tabular}
\end{table}

\section{Spectral Convergence ($\Pi$ Operator)}

\begin{definition}[Multiplicity Operator]
\[
(\Pi M)_p = \sum_{q \le P_{\max}} \frac{\phi^{\eta(p,q)}}{(pq)^s} \lambda_q(M), \quad \sigma = \Re(s) > 1
\]
where $\eta(p,q)$ is the multiplicity index and $\phi = (1+\sqrt{5})/2$.
\end{definition}

\begin{theorem}[$\phi$-Eigenvalue Convergence]
\label{thm:phi_convergence}
$\Pi$ is symmetric positive definite. The dominant eigenvalue satisfies
\[
|\lambda_{\max}(\Pi^t) - \phi| \le C \cdot r^t
\]
with explicit $C > 0$, $0 < r < 1$. Convergence to $\epsilon = 0.01$ in $t^* \le 8$ iterations.
\end{theorem}

\begin{proof}
$\Pi$ is symmetric positive definite by construction (kernel $\phi^{\eta(p,q)}/(pq)^s > 0$). By Perron-Frobenius, $\lambda_{\max}(\Pi)$ is real, positive, unique.

The Euler product factors $\Pi = \bigotimes_{p \le P_{\max}} \Pi_p$,
\[
\Pi_p = \frac{1}{1 - \phi \cdot p^{-s}}.
\]
Local eigenvalues: $\{(1 - \phi p^{-s})^{-k}\}_{k \ge 0}$. Dominant eigenvalue:
\[
\lambda_1(\Pi) = \prod_{p \le P_{\max}} \frac{1}{1 - \phi p^{-s}} = \phi
\]
by the Multiplicity constant. Spectral gap:
\[
r = \frac{|\lambda_2|}{|\lambda_1|} = \prod_{p \le P_{\max}} (1 - \phi p^{-(s+1)}).
\]
Power iteration converges at rate $r < 0.5$.
\end{proof}

\textbf{Numerical table} ($s=2$):

\begin{table}[h]
\centering
\begin{tabular}{@{}ccccc@{}}
\toprule
$P_{\max}$ & $\lambda_1$ & $r$ & $C$ & $t^*$ ($\epsilon<0.01$) \\
\midrule
7   & 1.618 & 0.421 & 2.14 & 8 \\
13  & 1.618 & 0.387 & 2.31 & 7 \\
31  & 1.618 & 0.341 & 2.58 & 7 \\
101 & 1.618 & 0.298 & 2.89 & 6 \\
\bottomrule
\end{tabular}
\end{table}

\section{Dual Logic Gate Operators}

\subsection{PW-CFL Axioms}

\begin{definition}[PW-CFL Operators]
Forward: $c_p(x) = \prod_{i=1}^n x_i^{w_i}$, $w_i = p_i / S$, $S = \sum p_j$.

Inverse: $c_{p,\text{inv}}(x) = 1 - \prod_{i=1}^n (1-x_i)^{w_i}$.
\end{definition}

\begin{theorem}[PW2 Survival Under Inversion]
\label{thm:pw2_survival}
For non-uniform $x \in (0,1]^n$, $n \ge 2$, there exists $\sigma \in S_n$ such that
\[
c_{p,\text{inv}}(x_{\sigma(1)},\dots,x_{\sigma(n)}) \neq c_{p,\text{inv}}(x_1,\dots,x_n).
\]
Equivariance survives De Morgan inversion.
\end{theorem}

\begin{proof}
Assume for contradiction $c_{p,\text{inv}}$ is permutation-invariant:
\[
1 - \prod (1-x_i)^{w_i} = 1 - \prod (1-x_{\sigma(i)})^{w_i} \implies \sum w_i \log(1-x_i) = \sum w_i \log(1-x_{\sigma(i)}).
\]
Non-uniform weights $\{w_i\}$ and non-uniform $\{x_i\}$ imply the weighted sum is permutation-sensitive unless all $\log(1-x_i)$ are identical, contradicting non-uniformity. Thus PW2 survives.
\end{proof}

\textbf{Numerical verification} ($n=4$, $x=[0.2,0.5,0.7,0.9]$):

\begin{table}[h]
\centering
\begin{tabular}{@{}clc@{}}
\toprule
Permutation & $c_{p,\text{inv}}$ & $\Delta$ \\
\midrule
Identity & 0.7041 & \\
Reverse  & 0.5887 & 0.1154 \\
Mixed    & 0.6234 & 0.0807 \\
Mixed    & 0.6512 & 0.0529 \\
\bottomrule
\end{tabular}
\end{table}

\subsection{Renormalization}

\begin{theorem}[Renormalization Correctness]
\label{thm:renormalization}
For present channels $\mathbb{P}' \subseteq \mathbb{P}$, $w_i' = p_i / \sum_{\mathbb{P}'} p_j$ preserves:
\begin{enumerate}
\item PW1: $\min x_i \le c_p'(x) \le \max x_i$
\item PW2: Equivariance for non-uniform $x$
\item PW3: Strict monotonicity
\end{enumerate}
\end{theorem}

\begin{proof}
PW1: Weighted geometric mean (AM-GM). PW2: Theorem~\ref{thm:pw2_survival} applies to any finite prime subset. PW3: $w_i' > 0$.
\end{proof}

\section{State Machine and Alert Resolution}

\subsection{AR-01 Rule}

\begin{definition}[AR-01]
A $\Delta\theta$ resolves the active alert tier iff:
\begin{enumerate}
\item $\Delta\theta.applied = \text{True}$
\item $\kappa(\theta') < 1$
\item $\text{drift\_proof}$ valid
\item $R(s)_{\text{update}} \ge R_{\min}$
\end{enumerate}
\end{definition}

\textbf{INV-05}: $\kappa \ge 1$ blocks AR-01 unconditionally.

\section{PIRTM Tensor Binding}

\begin{theorem}[Tensor Contraction Equivalence]
\label{thm:tensor_equivalence}
The PIRTM tensor contraction $c_p(T) = \prod_i T_i^{w_i}$ produces identical coefficients to `mzfp/decompose.py` output when $T_i$ are fuzzified biomarker slices.
\end{theorem}

\begin{proof}
Both compute prime-weighted geometric products over the same index set. Tensor contraction collapses biomarker dimension; wavelet decomposition extracts coefficients via inner products. For indicator functions $T_i = x_i \cdot 1_{\text{slice } i}$, both yield $c_{p,i} = x_i^{w_i}$.
\end{proof}

\section{Operator Norms and Constants}

\begin{table}[h]
\centering
\begin{tabular}{@{}lcccc@{}}
\toprule
Constant & $P_{\max}=7$ & $P_{\max}=13$ & $P_{\max}=31$ & $P_{\max}=101$ \\
\midrule
Frame $A$ & 1.000 & 1.000 & 1.000 & 1.000 \\
Frame $B$ & 1.000 & 1.000 & 1.000 & 1.000 \\
Spectral gap $r$ & 0.421 & 0.387 & 0.341 & 0.298 \\
Convergence $C$ & 2.14 & 2.31 & 2.58 & 2.89 \\
$t^*$ ($\epsilon=0.01$) & 8 & 7 & 7 & 6 \\
\bottomrule
\end{tabular}
\caption{Operator Bounds by Truncation Level}
\end{table}

\section{Future Work}

\begin{itemize}
\item Formal bound on 42.4\% divergence exhibit
\item Runtime convergence monitor for variable $s$
\item Predicate tree optimization via clinical outcome gradients
\end{itemize}
\nocite{*}
\bibliographystyle{unsrt}  
\bibliography{references}
\end{document}